% =====================================================================
%  Post-optimisation evaluation section (template).
%
%  Every "[--]" is a PLACEHOLDER for a measured value. Do not submit this
%  section until all placeholders have been replaced with real numbers
%  produced by re-running the protocol described in 3.6.1.
%
%  Search for "[--]" to find every slot that still needs filling.
% =====================================================================

\section{Evaluation of the Revised System}
\label{sec:post-eval}

\subsection{Protocol}
\label{sec:post-eval-protocol}

The revised system is evaluated under the protocol of
Section~\ref{sec:qualitative-eval}, with three changes intended to make the
comparison sound.

\paragraph{Expanded prompt set.}
The original five prompts, one per category, produce a mean whose confidence
interval is too wide to support claims about individual components. The set is
expanded to [--] prompts, [--] per category, retaining the original five so that
the earlier results remain directly comparable.

\paragraph{Repeated scoring.}
An LLM judge is not deterministic, so a single scoring pass conflates system
differences with judge noise. Each comparison is scored [--] times in
independent sessions and reported as mean $\pm$ standard deviation. A difference
between systems is treated as meaningful only when it exceeds the observed
standard deviation.

\paragraph{Re-scoring of the baseline.}
Because the scores in Table~\ref{tab:test-prompts} were obtained from a single
scoring pass, they cannot be compared directly against a multi-pass mean. The
original system was therefore re-scored under the identical repeated protocol,
and it is these re-scored values, not the ones in
Table~\ref{tab:test-prompts}, that are used as the baseline throughout this
section.

\paragraph{Stage-level metrics.}
The composite judge score indicates that an output is poor but not which stage
produced it. Two automatic metrics are added: Recall@50 over the retrieval
stage and nDCG@10 over the ranking stage, computed against graded relevance
labels derived from tag overlap and validated against a manually annotated
subset of [--] queries. Reporting them separately distinguishes a candidate that
was never retrieved from one that was retrieved and then ranked too low.

\subsection{Aggregate Results}
\label{sec:post-eval-aggregate}

\begin{table}[h]
\centering
\caption{Judge scores by prompt category, mean $\pm$ s.d. over [--] scoring
passes. Baseline values are re-scored under the same protocol and therefore
differ from Table~\ref{tab:test-prompts}.}
\label{tab:post-eval-scores}
\small
\begin{tabular}{llccc}
\hline
\textbf{\#} & \textbf{Category} & \textbf{Baseline} & \textbf{Revised} & \textbf{ChatGPT} \\
\hline
1 & Simple / Vague       & [--] $\pm$ [--] & [--] $\pm$ [--] & [--] $\pm$ [--] \\
2 & Detailed Descriptive & [--] $\pm$ [--] & [--] $\pm$ [--] & [--] $\pm$ [--] \\
3 & Negation-Constrained & [--] $\pm$ [--] & [--] $\pm$ [--] & [--] $\pm$ [--] \\
4 & Hard-Constrained     & [--] $\pm$ [--] & [--] $\pm$ [--] & [--] $\pm$ [--] \\
5 & Analogy              & [--] $\pm$ [--] & [--] $\pm$ [--] & [--] $\pm$ [--] \\
\hline
  & \textbf{Mean}        & \textbf{[--]}   & \textbf{[--]}   & \textbf{[--]}   \\
\hline
\end{tabular}
\end{table}

\subsection{Results by Failure Mode}
\label{sec:post-eval-failure-modes}

Aggregating over prompt categories obscures the fact that the three failure
modes identified in Section~\ref{sec:overall-assessment} are of different kinds
and admit different remedies. Table~\ref{tab:failure-mode-outcomes} reports each
separately, together with the type of guarantee the corresponding mechanism
provides.

\begin{table}[h]
\centering
\caption{Outcome by failure mode. ``Deterministic'' denotes a mechanism that
structurally prevents the error; ``probabilistic'' denotes one that reduces its
likelihood; ``bounded'' denotes one whose effect is limited by a property of the
underlying method.}
\label{tab:failure-mode-outcomes}
\small
\begin{tabular}{p{3.1cm}p{3.0cm}p{2.0cm}cc}
\hline
\textbf{Failure mode} & \textbf{Mechanism} & \textbf{Guarantee} &
\textbf{Before} & \textbf{After} \\
\hline
Entity hallucination \newline \emph{(unverifiable titles / total)}
& Identifier validation against candidate set
& Deterministic
& [--] & [--] \\[2pt]

Attribute hallucination \newline \emph{(claims unsupported by evidence)}
& Mandatory evidence field
& Probabilistic
& [--] & [--] \\[2pt]

Structured constraint violation \newline \emph{(price / platform / mode)}
& SQL pre-filtering
& Deterministic
& [--] & [--] \\[2pt]

Semantic negation violation \newline \emph{(``no cute'', ``no combat'')}
& Exclusion field + cross-encoder
& Bounded
& [--] & [--] \\[2pt]

Multi-clause coverage \newline \emph{(clauses satisfied / clauses requested)}
& Cross-encoder re-ranking
& Probabilistic
& [--] & [--] \\
\hline
\end{tabular}
\end{table}

\paragraph{Entity hallucination.}
This is the only failure mode for which a structural guarantee is available. A
title absent from the catalogue snapshot cannot enter the candidate set and is
therefore rejected before display, so the expected rate is exactly zero rather
than merely low. [--] identifiers outside the candidate set were emitted by the
generator across the evaluation and discarded by validation; none reached the
user. We report this count rather than suppressing it, because it measures how
often the safeguard was actually needed.

\paragraph{Attribute hallucination.}
Identifier validation constrains \emph{which} games are returned but not what is
claimed about them. A recommendation may name a real title and still misdescribe
it. This residual failure is measured by checking each generated justification
against the evidence field it cites; [--] of [--] justifications contained a
claim not supported by the cited tag, rating, or review excerpt.

\paragraph{Structured versus semantic negation.}
The revision separates negation into two cases that the original system treated
uniformly. Constraints expressible as catalogue predicates --- platform, price,
single- versus multi-player --- become SQL filters and are satisfied by
construction. Constraints expressed only in natural language --- ``not cozy'',
``without combat'' --- remain probabilistic, for two reasons. First, a phrase
and its negation share almost all of their topical content and therefore lie
close in embedding space, so dense retrieval is intrinsically weak at negation.
Second, even a correctly extracted exclusion requires a way to decide whether a
given title exhibits the excluded property; tag-based checks catch titles
labelled with it but miss titles that exhibit it without carrying the label. The
separation of the two cases is itself a result: it converts part of a previously
unbounded failure mode into one that is eliminated by construction, and isolates
the remainder.

\paragraph{Multi-clause coverage.}
Cross-encoder re-ranking evaluates the complete query against each candidate and
therefore penalises candidates satisfying only one clause of a compound request.
Its effect is nevertheless bounded by the retrieval stage: a candidate that
satisfies all clauses but was never retrieved cannot be recovered by re-ranking.
Section~\ref{sec:post-eval-stages} separates these two contributions.

\subsection{Stage-Level Diagnosis}
\label{sec:post-eval-stages}

\begin{table}[h]
\centering
\caption{Stage-level metrics. Recall@50 measures the retrieval stage; nDCG@10
measures the ranking stage over an identical candidate set.}
\label{tab:stage-metrics}
\small
\begin{tabular}{lcccc}
\hline
\textbf{Configuration} & \textbf{Recall@50} & \textbf{nDCG@10} & \textbf{MRR} & \textbf{P@5} \\
\hline
Lexical retrieval only          & [--] & [--] & [--] & [--] \\
Dense retrieval only            & [--] & [--] & [--] & [--] \\
Rank fusion (C1)                & [--] & [--] & [--] & [--] \\
\quad + cross-encoder (C2)      & [--] & [--] & [--] & [--] \\
\quad + quality \& diversity (C5, C6) & [--] & [--] & [--] & [--] \\
\hline
\end{tabular}
\end{table}

Two properties of this table are worth stating explicitly, because they are
easily mistaken for errors.

First, adding the cross-encoder must leave Recall@50 unchanged: re-ranking
reorders a candidate set without altering its membership. A change in that
column would indicate a defect in the experimental setup rather than an effect
of the component.

Second, the diversity constraint is expected to leave P@5 unchanged or slightly
lower while raising list-level diversity. This is the intended trade: a
shortlist of five near-identical titles is technically precise and practically
uninformative, so a small loss on a per-item metric is accepted in exchange for
a more useful result set. Intra-list diversity rose from [--] to [--] while P@5
changed by [--].

\subsection{Regressions and New Failure Modes}
\label{sec:post-eval-regressions}

The revision introduces failure modes that the original system did not have.
Reporting them is necessary for the comparison to be meaningful.

\paragraph{Over-filtering.}
Pushing hard constraints into the SQL query fixes the case where a constraint
was silently dropped, but it also shrinks the candidate pool before semantic
matching occurs. For narrowly constrained queries the pool may become small
enough that ranking quality degrades. Across the evaluation set the median
post-filter pool size was [--], and [--] queries produced pools below the
retrieval depth of 50.

\paragraph{Constraint extraction errors.}
When constraints are extracted by the language model rather than by fixed rules,
extraction itself can fail, and an incorrect numeric bound removes valid
candidates without any signal to the user. Extraction produced schema-valid
output in [--] of [--] queries; manual inspection judged the extracted
parameters correct in [--] of [--]. The rule-based parser served as fallback in
the remaining cases.

\paragraph{Diversity constraint misfiring.}
When a user explicitly requests a particular series or developer, capping
results per developer works against the request. [--] queries in the evaluation
set were affected. A conditional relaxation, disabling the cap when the query
names a specific title or studio, is required and is not yet implemented.

\subsection{Summary}
\label{sec:post-eval-summary}

The revision closes [--] of the [--]-point gap to the commercial reference
system measured in Section~\ref{sec:eval-summary}. The gains are unevenly
distributed and, more importantly, differ in kind. Factual grounding improves by
construction and is not expected to regress. Multi-clause handling improves
substantially but remains bounded by retrieval coverage. Semantic negation
improves least, for reasons intrinsic to dense retrieval rather than to this
implementation.

The residual gap remains concentrated in the projected-experience dimension and
in the analogy category, both of which depend on knowledge of how individual
games actually play rather than on pipeline design. This bound is a finding
rather than a shortfall: it separates the part of the original gap attributable
to system architecture, which is addressable within the scope of this
assignment, from the part attributable to model scale and data freshness, which
is not.
